% =====================================================================
% APPENDIX — robustness to the quality threshold and to the sigma_1 = 0
% restriction. Referenced from §4.1.
% =====================================================================

\section{Robustness: The Quality Threshold and the $\sigma_1^*=0$ Restriction}
\label{app:robustness}

Two maintained choices in the baseline specification deserve separate scrutiny, because each
is an assumption rather than an estimate.

The first is where the binary quality space is cut. The baseline places the boundary between
D and E, so that $\theta=1$ corresponds to grades A--D. This follows the rental-ban
discontinuity of \emph{loi Climat et R\'esilience} discussed in
Section~\ref{ss:institutions}, and it is the cut for which the price-index contrast and the
listing-level quality measure refer to the same partition of grades. The natural alternative
places the boundary between C and D.

The second is the treatment of informed sellers who withhold. Because such a listing displays
no score, its quality cannot be read from the listing itself. The baseline assigns
$\theta=0$ by revealed preference, invoking Corollary~\ref{cor:hq_disclose}. That assignment
is an assumption, and it has a consequence worth stating plainly: it makes
$\hat\sigma_{1,m}$ identically zero by construction, so the corollary cannot be assessed
against the baseline estimates. The alternative reads the withheld grade from the ADEME
registry, which records it for every matched listing, leaving $\hat\sigma_{1,m}$ measured.

Table~\ref{tab:robustness} crosses the two choices.

\begin{table}[H]
\centering
\caption{Sensitivity to the quality threshold and to the $\sigma_1$ assignment}
\label{tab:robustness}
\small\renewcommand{\arraystretch}{1.25}
\begin{tabular*}{\textwidth}{@{\extracolsep{\fill}}l*{6}{c}@{}}
\toprule\toprule
\textbf{Specification} & $N$ & $\hat\rho_m$ & $\hat c_{0,m}$ & $\hat\sigma_{0,m}$ & $\hat\sigma_{1,m}$ & $\hat{\bar\theta}_m$ \\
\midrule
\multicolumn{7}{l}{\textit{Panel A: Apartments}} \\[2pt]
A--D vs E/F/G, $\sigma_1$ imposed \textbf{(baseline)} & 45 & 90.9\% & 3.42\% & 14.0\% & --- & 62.5\% \\
A--D vs E/F/G, $\sigma_1$ from registry               & 45 & 90.8\% & 3.75\% &  5.5\% & 5.1\% & 65.9\% \\
A--C vs E/F/G, $\sigma_1$ imposed                     & 40 & 95.0\% & 1.62\% &  7.1\% & --- & 26.5\% \\
A--C vs E/F/G, $\sigma_1$ from registry               & 40 & 95.0\% & 1.75\% &  4.8\% & 6.5\% & 28.4\% \\
\midrule
\multicolumn{7}{l}{\textit{Panel B: Houses}} \\[2pt]
A--D vs E/F/G, $\sigma_1$ imposed \textbf{(baseline)} & 97 & 68.1\% & 4.45\% &  5.7\% & --- & 60.5\% \\
A--D vs E/F/G, $\sigma_1$ from registry               & 97 & 67.7\% & 4.59\% &  2.2\% & 2.1\% & 61.9\% \\
A--C vs E/F/G, $\sigma_1$ imposed                     & 100 & 71.8\% & 1.80\% &  3.0\% & --- & 28.5\% \\
A--C vs E/F/G, $\sigma_1$ from registry               & 100 & 71.7\% & 1.88\% &  1.9\% & 2.5\% & 29.3\% \\
\bottomrule\bottomrule
\end{tabular*}
\vspace{0.5em}
\begin{minipage}{\linewidth}
\footnotesize\textit{Notes.} Market-level means. ``$\sigma_1$ imposed'' assigns $\theta=0$ to
informed non-disclosers by revealed preference, which forces $\hat\sigma_{1,m}=0$ in every
market; the column is left blank there because the quantity is assumed rather than estimated.
``$\sigma_1$ from registry'' binarises the withheld grade from the ADEME record. The quality
threshold also governs the price-index contrast: A--D vs E/F/G merges the \emph{Efficient} and
\emph{Balanced} index series by transaction volume, while A--C vs E/F/G contrasts
\emph{Efficient} against \emph{Intensive} and omits \emph{Balanced}. Market counts differ
across rows because the minimum-volume filter is applied to the high-group series, which is
larger under the A--D contrast.
\end{minipage}
\end{table}

Three results follow.

First, the structural estimates are insensitive to the $\sigma_1$ assignment. Holding the
threshold fixed, relaxing the restriction moves $\hat\rho_m$ by less than a tenth of a
percentage point in both segments, and moves $\hat c_{0,m}$ from 3.42\% to 3.75\% for
apartments and from 4.45\% to 4.59\% for houses. The concealment cost is therefore not an
artefact of the revealed-preference assignment, even though $\hat\sigma_{0,m}$ itself is
sensitive to it: imposing $\theta=0$ places every informed non-discloser in the low-quality
group and roughly doubles the measured hiding rate.

Second, the restriction $\sigma_1^*=0$ is not satisfied exactly, but the model's predicted
\emph{ordering} holds under the maintained threshold. With quality read from the registry and
the cut at the rental-ban boundary, $\hat\sigma_{1,m} = 5.1\%$ against
$\hat\sigma_{0,m} = 5.5\%$ for apartments, and $2.1\%$ against $2.2\%$ for houses: informed
sellers of low-quality properties withhold more often than informed sellers of high-quality
ones, as the model requires, though the gap is narrow. Under the C--D cut the ordering
reverses in both segments ($6.5\%$ against $4.8\%$, and $2.5\%$ against $1.9\%$). The
threshold that the institutional setting selects is thus also the one under which the
behavioural restriction survives, which we read as support for the baseline cut rather than
as independent confirmation of the corollary.

Third, the threshold matters a great deal for the level of the concealment cost, and this is
the paper's most consequential specification choice. Moving the cut from C--D to D--E raises
$\hat c_{0,m}$ from roughly 1.6--1.8\% of property value to 3.4--4.6\%. The mechanism is
transparent: with A--D counted as high quality, the average quality of the non-disclosing
pool $\hat e^*_m$ roughly doubles, and $c_{0,m} = \rho_m\beta_{1,m}e^*_m$ scales with it.
Readers who prefer the narrower definition of ``high quality'' should therefore read the
concealment cost as the smaller figure throughout; the qualitative results --- the ordering
across segments, the boundary saturation in $\hat\rho_m$, and the regime classification ---
are unaffected.
